% EDITION STATUS: partial; source units 3/18
\documentclass[11pt,a4paper]{article}
\usepackage[no-math]{fontspec}
\setmainfont{Latin Modern Roman}
\setsansfont{Latin Modern Sans}
\setmonofont{Latin Modern Mono}
\usepackage[french,english]{babel}
\usepackage[left=22mm,right=34mm,top=22mm,bottom=22mm,marginparwidth=23mm,marginparsep=5mm,includeheadfoot]{geometry}
\usepackage{amsmath,amssymb}
\usepackage[normalem]{ulem}
\usepackage{xcolor,tikz-cd,marginnote,etoolbox,microtype}
\usepackage[unicode,colorlinks=true,linkcolor=LinkInk,urlcolor=LinkInk,bookmarksopen=true]{hyperref}
\definecolor{LinkInk}{HTML}{254E70}
\definecolor{HandInk}{HTML}{254E70}
\definecolor{DeletionInk}{HTML}{8B3A3A}
\definecolor{QuietInk}{HTML}{555555}
\setlength{\parindent}{1em}
\setlength{\parskip}{3pt plus 1pt}
\setlength{\emergencystretch}{1.5em}
\raggedbottom
\normalmarginpar
\frenchspacing
\newif\ifdiplomatic
\newif\ifreading
\newif\ifenglish
\diplomaticfalse
\readingfalse
\englishtrue
\newcommand{\EditionTitle}{A Guide to the Esquisse}
\newcommand{\EditionForm}{Commentary for a mathematical reader}
\hypersetup{pdftitle={Sketch of a Program — Corrected text with translator’s notes},pdfauthor={Alexander Grothendieck},pdfsubject={A scan-checked edition of Esquisse d'un programme, January 1984, Montpellier cote 119}}
\newcommand{\DeclareScan}[4]{%
  \expandafter\def\csname edurl@#1\endcsname{#2}%
  \expandafter\def\csname edkey@#1\endcsname{#3}%
  \expandafter\def\csname edpage@#1\endcsname{#4}}
\def\CurrentScanURL{}
\def\CurrentPageKey{}
\def\CurrentCopy{copy1}
\def\CurrentApparatus{}
\newcounter{appnote}
\newcommand{\EditionPart}[2]{\clearpage\def\CurrentCopy{#1}\section*{#2}}
\newcommand{\scanpage}[1]{%
  \clearpage
  \expandafter\let\expandafter\CurrentScanURL\csname edurl@#1\endcsname
  \expandafter\let\expandafter\CurrentPageKey\csname edkey@#1\endcsname
  \def\CurrentApparatus{}%
  \hypertarget{scan-\CurrentPageKey}{}%
  \typeout{EDITION-SCAN: \CurrentPageKey; PDFPAGE: \the\numexpr\ReadonlyShipoutCounter+1\relax}%
  \ifreading
    \marginnote{\raggedright\scriptsize\sffamily\href{\CurrentScanURL}{119 / \csname edpage@#1\endcsname}}%
  \else
    \noindent{\small\sffamily\href{\CurrentScanURL}{[scan #1]}}\par\smallskip
  \fi}
\newcommand{\mspage}[1]{%
  \ifreading\marginnote{\raggedright\scriptsize\sffamily p.\,#1}[2\baselineskip]%
  \else\noindent{\small\sffamily[p.\,#1]}\par\smallskip\fi}
\newcommand{\hand}[1]{\ifdiplomatic\textcolor{HandInk}{#1}\else#1\fi}
\let\over\relax
\newcommand{\over}[2]{\textcolor{QuietInk}{#1}\textcolor{HandInk}{\ensuremath{^{\!\nearrow}}#2}}
\newcommand{\inserted}[2][]{\textcolor{HandInk}{\ensuremath{\langle}#2\ensuremath{\rangle}}}
\newcommand{\deleted}[1]{\ifmmode\textcolor{DeletionInk}{[#1]_{\scriptscriptstyle\mathrm{del}}}\else\textcolor{DeletionInk}{\sout{#1}}\fi}
\renewcommand{\xout}[1]{\ifmmode\textcolor{QuietInk}{[#1]_{\scriptscriptstyle\mathrm{x}}}\else\textcolor{QuietInk}{\sout{#1}}\fi}
\let\TextUline\uline
\let\TextUuline\uuline
\renewcommand{\uline}[1]{\ifmmode\underline{#1}\else\TextUline{#1}\fi}
\renewcommand{\uuline}[1]{\ifmmode\underline{\underline{#1}}\else\TextUuline{#1}\fi}
\newcommand{\unsure}[1]{\href{\CurrentScanURL}{[#1?]}}
\newcommand{\illegible}[1]{\href{\CurrentScanURL}{[\ifenglish illegible\else illisible\fi\ #1]}}
\newsavebox{\MarginalTestBox}
\newcommand{\margin}[1]{%
  \sbox{\MarginalTestBox}{\parbox{\marginparwidth}{\raggedright\footnotesize#1}}%
  \ifdim\dimexpr\ht\MarginalTestBox+\dp\MarginalTestBox\relax>3\baselineskip
    \par\smallskip\noindent\begin{minipage}{\linewidth}\hrule\smallskip
    {\footnotesize\textcolor{HandInk}{\ifenglish Marginal note:\else Note marginale :\fi\ #1}}\smallskip\hrule\end{minipage}\par\smallskip
  \else
    \marginnote{\raggedright\footnotesize\textcolor{HandInk}{#1}}%
  \fi}
\newcommand{\authornote}[1]{\margin{#1}}
\newcommand{\diagram}[1]{\par\smallskip{\centering#1\par}\smallskip}
\newenvironment{diagramblock}{\par\begin{center}}{\end{center}\par}
\newcommand{\DeclareVariant}[2]{\expandafter\def\csname edvariant@#1\endcsname{#2}}
\newcommand{\AppItem}[2]{\par\noindent\textsuperscript{#1}\csname edvariant@#2\endcsname}
\newcommand{\EditionApp}[1]{%
  \stepcounter{appnote}%
  \ifmmode{}^{\scriptscriptstyle\theappnote}\else\textsuperscript{\theappnote}\fi
  \begingroup\edef\x{\endgroup\noexpand\gappto\noexpand\CurrentApparatus{\noexpand\AppItem{\theappnote}{#1}}}\x}
\newcommand{\var}[4]{\ifdiplomatic#1\else#2\ifreading\else\EditionApp{#4}\fi\fi}
\newcommand{\copytwo}[2]{\ifdiplomatic\else\ifreading\else\DeclareVariant{#1}{#2}\EditionApp{#1}\fi\fi}
\newcommand{\tnote}[2]{\ifreading\else\DeclareVariant{#1}{\textit{Translator's note.} #2}\EditionApp{#1}\fi}
\newcommand{\finishscan}{\ifx\CurrentApparatus\empty\else\par\bigskip\noindent\rule{.35\linewidth}{.35pt}\par\smallskip{\footnotesize\setlength{\parskip}{3pt}\CurrentApparatus}\fi}
\newcommand{\DeclareNoteReturns}[2]{\expandafter\def\csname edreturns@#1\endcsname{#2}}
\newcommand{\notecall}[3]{\hypertarget{#2}{\hyperlink{note-#1}{#3}}}
\newcommand{\notedef}[2]{\hypertarget{note-#1}{#2}\,{\scriptsize[\csname edreturns@#1\endcsname]}}
\newcommand{\AppReading}[1]{{\def\margin##1{##1}\def\authornote##1{##1}#1}}
\newcommand{\EmptyReading}{\ifenglish\textit{omitted}\else\textit{supprimé}\fi}



\begin{document}
\begin{titlepage}
\centering
\vspace*{18mm}
{\Large Alexander Grothendieck\par}
\vspace{12mm}
{\Huge\bfseries\EditionTitle\par}
\vspace{8mm}
{\Large\EditionForm\par}
\vspace{14mm}
{\large January 1984\par}
\vfill
{\large Université de Montpellier\par fonds Grothendieck, cote 119\par}
\vspace{8mm}
{\normalsize Two typescript copies: sites 2--59 and 78--135\par}
\vspace{8mm}
{\normalsize Edition dated 14 November 2026\par}
\end{titlepage}
\thispagestyle{plain}
\section*{Preface}
\noindent\textbf{Working checkpoint.} This guide currently contains 3 of the 18 requested commentary units. Its scope is recorded with the checked-page coverage.\par\smallskip
This guide follows the sommaire, sections 1–10 and notes (1)–(7) in the order of the typescript. Page references give the original typescript pagination and link to the archive. The guide is based on the checked text and collation; it does not change the text’s mathematics.

\smallskip
AI-transcribed by Claude Opus 5.5 (Anthropic, effort high, Message Batches API); checked against the scans, and this edition built, by GPT-6 Astra (OpenAI, ChatGPT Pro mode); prompts and records by Claude (Anthropic).


The transcription of copy 1 was made with the text layer of the Grothendieck Circle's typeset \textit{Esquisse} (\texttt{EsquisseFr.pdf}) given to the transcribing model as a draft, to compare with after it had transcribed each page; every reading of this edition was decided on the scans.

The published French text and English translation in \textit{Geometric Galois Actions 1}, edited by L. Schneps and P. Lochak (Cambridge University Press, 1997), the Grothendieck Circle's typeset \texttt{EsquisseFr.pdf}, and the Centre for Grothendieckian Studies' draft transcription G94d are named as published texts. They are not reproduced here.

\section*{Rights and credits}
The rights in the original work belong to its author or the author's representatives or heirs. Where we hold rights, in the AI-produced transcription, translation, corrections and typesetting, they are dedicated to the public domain under CC0 1.0. This edition is published for educational, non-profit purposes, and we claim no rights we do not hold. If you hold rights in this work and object to its publication here, please open a GitHub issue and we will take it down.

\clearpage


\subsection*{Sommaire}
The contents page arranges the programme in ten parts: the \emph{Envoi}; the
``jeu de Légo-Teichmüller'' and the Galois group; number fields attached to a
``dessin d'enfant''; regular polyhedra and finite fields; the passages toward
``topologie modérée'' and ``théories modérées''; ``A la Poursuite des Champs'';
two-dimensional geometry; teaching; and the epilogue. This ordering gives the
reader two practical routes into the mathematical discussion, through the
surfaces and combinatorics of \S\S2--4 and through the proposed foundations of
\S\S5--7, before the return to two-dimensional geometry in \S8. The N.B.
also specifies the document's two systems of notes: asterisks refer to notes
at the bottom of the same page, while the numbered calls refer to the notes
``rajoutées ultérieurement'' collected at the end. In copy 1 the endpoints
1 and 7 of that numbered sequence are handwritten above typed parentheses;
the small dot beside 7 remains marked as doubtful. The French forms retain
the scanned spellings ``astéristiques'' and ``bidimensionelle'', as well as
the accent in ``Légo''. These details belong to the reader's route through
the document and its successive additions.
(\href{https://grothendieck.umontpellier.fr/web/viewer.html?file=../119.pdf\#page=3}{Copy 1, sommaire, site 2}.)



\subsection*{1. Envoi}
The opening locates the programme within Grothendieck's application to the
CNRS and explains the work he intends to carry forward. He refers to the
1972 account of his ``Titres et Travaux'', with its list of publications,
and to two boxes of subsequent manuscript notes that have passed through
successive stages of ``décantation''. His account of ``rédaction'' already
qualifies the word: on p.~1 he describes the projected writing as the
further development, clarification, and deepening of ideas, with the
unexpected turns that accompany discovery. The ``Réflexions Mathématiques'',
begun in February 1983 according to this passage, are to make that movement
of thought visible. The sentence crossing pp.~1--2 gives the complementary
operation: an inadequate image is exposed, and the objects previously
encountered askew acquire their mutual coherence. The final paragraph
assigns the ensuing sketch both a retrospective role, covering ten or
twelve years of reflection, and a prospective role as a programme for the
coming years. Its stated readers include the Comité National and other
colleagues, former students, and friends. The 1972 report and the manuscript
``Réflexions Mathématiques'' are named here as the author's own points of
reference; the checked passage supplies their role in this application.
(\href{https://grothendieck.umontpellier.fr/web/viewer.html?file=../119.pdf\#page=4}{Copy 1, p.~1, site 3};
\href{https://grothendieck.umontpellier.fr/web/viewer.html?file=../119.pdf\#page=5}{p.~2, site 4}.)



\subsection*{2. A ``jeu de Légo-Teichmüller'' and the Galois group}
The second section develops the programme from the shared geometric
intuition of surfaces that Grothendieck describes as emerging from his
teaching (pp.~2--3). The transition from combinatorial topology to complex
algebraic curves brings back the fields, rings, and groups named on p.~3,
then the ``multiplicités modulaires'' and their automorphisms. On
pp.~4--5 the family of moduli multiplicities, its gluing and erasing
operations, and the corresponding fundamental groupoids form the
``tour de Teichmüller'', on which he considers the action of the absolute
Galois group. He states the faithfulness of the action on the first
nontrivial level and asks for a description of the tower's action in terms
of that level; p.~5 explicitly says that he has no conjecture yet for a
characterization of the Galois group inside the outer automorphism group.
The groupoid language keeps the chosen collections of base points visible
under the symmetries. On pp.~6--7 he formulates the construction principle
in which the first two moduli dimensions supply generators and relations,
sets it beside Demazure's rank principle, and describes the remaining
topological work in terms of the fundamental group at infinity and
``dévissage'' of stratified structures, with an explicit reference to
\S5. On pp.~8--9 the ``petite boite'' and ``grande boite'' turn those
operations into assemblies of surfaces; Thurston's surgery is named,
canonical cellular decompositions point forward to \S3, and real
structures guide the proposed decompositions of moduli multiplicities.
A mathematical slip remains printed on p.~9: the relations
$\rho^3=\sigma^2$, $\rho^4=\sigma^6=1$ imply that
$z=\rho^3=\sigma^2$ satisfies $z^4=z^3=1$, hence $z=1$,
$\rho=1$, and $\sigma^2=1$. The resulting presentation has mutually
inverse maps to the cyclic group of order two, sending $\rho$ to its
identity and $\sigma$ to its generator, and sending that generator back
to $\sigma$. The matrix $\left(\begin{smallmatrix}1&1\\0&1\end{smallmatrix}\right)$
has infinite order in $\mathrm{SL}_2(\mathbb Z)$, which establishes the
failure of the displayed presentation to present the group named there.
The readers retain the scanned exponents. For subsequent work on the
Galois-group programme, one may name \emph{La Longue Marche à travers la
théorie de Galois} and Drinfeld's
\href{https://www.mathnet.ru/eng/aa199}{\emph{On quasitriangular quasi-Hopf algebras and on a group that is closely connected with}
$\mathrm{Gal}(\overline{\mathbb Q}/\mathbb Q)$}.
(\href{https://grothendieck.umontpellier.fr/web/viewer.html?file=../119.pdf\#page=5}{Copy 1, pp.~2--9, beginning at site 4};
\href{https://grothendieck.umontpellier.fr/web/viewer.html?file=../119.pdf\#page=12}{mathematical slip on p.~9, site 11}.)
% Complete the explicit backward reference to note (1) when its scan check is finished.


\end{document}
